\section{Data Collection}
\label{sec:data}

This section describes the data collection pipeline for yield-bearing token activity and whale identification across Ethereum mainnet lending protocols.

\subsection{Token Selection}

Three USD-denominated yield-bearing tokens were selected based on the following criteria: top candidates by Ethereum mainnet circulating value and TVL, minimum on-chain liquidity threshold, minimum active holder count, active integration as collateral on at least one major lending protocol, and activity throughout the sample period.

\begin{table}[h]
\centering
\caption{Selected yield-bearing tokens.}
\label{tab:tokens}
\begin{tabular}{llll}
\hline
\textbf{Token} & \textbf{Issuer} & \textbf{Market Cap} & \textbf{Typical APY} \\
\hline
sUSDe & Ethena & \$3.5B & 4--8\% \\
sUSDS & Sky & \$6B & 4--6\% \\
syrupUSDC & Maple Finance & \$1.8B & 5--8\% \\
\hline
\end{tabular}
\end{table}

These tokens represent the largest USD-denominated yield-bearing stablecoins on Ethereum mainnet with active lending protocol integrations.

\subsection{Protocol Coverage}

Three lending protocols are covered based on where these tokens are actively used as collateral.

\begin{table}[h]
\centering
\caption{Protocol coverage by token.}
\label{tab:protocols}
\begin{tabular}{lcc}
\hline
\textbf{Token} & \textbf{Aave v3} & \textbf{Morpho Blue} \\
\hline
sUSDe & Yes & Yes \\
sUSDS & No & Yes \\
syrupUSDC & No & Yes \\
\hline
\end{tabular}
\end{table}

\subsection{YBT Scraper}

The YBT scraper collects metadata, prices, and protocol events for the three tokens under study. Data is sourced from DeFiLlama (TVL, APY) and Dune Analytics (prices, events).

\begin{table}[h!]
\centering
\caption{YBT scraper output files.}
\label{tab:ybt-output}
\begin{tabular}{lp{7cm}}
\hline
\textbf{Output File} & \textbf{Contents} \\
\hline
\texttt{ybt\_metadata.json} & Token name, issuer, address, decimals, TVL, APY \\
\texttt{ybt\_prices.csv} & Daily average, low, and high prices for depeg monitoring \\
\texttt{aave\_events.csv} & Supply, withdraw, borrow, repay, liquidation events (sUSDe) \\
\texttt{morpho\_events.csv} & Events for all three tokens on Morpho Blue \\
\hline
\end{tabular}
\end{table}

Each event record includes timestamp, block number, transaction hash, user address, and amount.

\subsection{Whale Identification}

The whale identification pipeline finds the top addresses per token per protocol that have actively engaged in looping or flashloan strategies, ranked by lifetime token volume.

\textbf{Pipeline:}
\begin{enumerate}
    \item Run Dune SQL queries per token per protocol, joining ERC-20 Transfer events with protocol-specific supply, borrow, and flashloan events.
    \item Flag each address as EOA or smart contract via \texttt{eth\_getCode} RPC check.
    \item Output ranked whale lists per token as CSV.
\end{enumerate}

\textbf{Strategy definitions:}
\begin{itemize}
    \item \textbf{Looper}: supplied the token as collateral at least once and borrowed any asset at least 3 times on the same protocol.
    \item \textbf{Flashloaner}: executed flashloans in at least 3 separate transactions on the respective protocol.
\end{itemize}

\textbf{Fields collected per whale:}
\begin{itemize}
    \item \texttt{rank}: position by total received volume
    \item \texttt{address}: wallet address
    \item \texttt{current\_balance}: current token holdings
    \item \texttt{total\_received}: lifetime cumulative inflow
    \item \texttt{distinct\_funders}: number of unique addresses that sent tokens to this wallet
    \item \texttt{is\_eoa}: True if externally owned account, False if smart contract
    \item \texttt{is\_looper}, \texttt{is\_flashloaner}: strategy flags
    \item \texttt{strategy\_type}: looper, flashloaner, or both
    \item \texttt{protocol}: Aave, Morpho
\end{itemize}

Infrastructure contracts such as the Morpho ETH Bundler are excluded as they aggregate transactions from many end users rather than representing independent strategy actors.


\subsection{Selected Whales}
After running the identification queries and manually reviewing each candidate on Etherscan, six addresses were selected for the core analysis. The selection criteria required: a single identifiable operator (not a multi-user vault), sufficient transaction history for panel analysis, and no contamination from unrelated strategies such as Pendle PT looping. Several additional candidates were rejected after closer inspection of their ERC-20 transfer history revealed mixed strategies or multi-user aggregation.

\begin{table}[h!]
\centering
\caption{Selected whales for analysis.}
\label{tab:whales}
\begin{tabular}{llllll}
\hline
\textbf{ID} & \textbf{Address} & \textbf{Protocol} & \textbf{Collateral} & \textbf{Borrow} & \textbf{Strategy} \\
\hline
W-1 & 0x080c...fb0 & Morpho & sUSDe & USDtb \\
W-2 & 0x97bb...21e & Morpho & sUSDS & USDT \\
W-3 & 0xcb1e...e58 & Morpho & sUSDS & USDT \\
W-4 & 0x0d4f...a06b & Morpho & sUSDe & PYUSD \\
W-5 & 0x2438...6717 & Morpho & syrupUSDC & PYUSD \\
W-6 & 0x0294...080 & Aave v3 & sUSDe & USDT \\
W-7 & 0x2420...1e3 & Aave v3 & sUSDe & USDe \\
\hline
\end{tabular}
\end{table}

The sample covers all three tokens across two protocols. W-4 and W-5 share the same deployer and operate as a coordinated parallel strategy across different collateral tokens. W-6 and W-7 both employ Aave v3 E-Mode looping on sUSDe, with W-7 distinguished by its intraday open-and-unwind cadence.
\subsection{Data Sources and Pipeline}
The complete data pipeline uses the following sources:
\begin{table}[h!]
\centering
\caption{Data sources.}
\label{tab:sources}
\begin{tabular}{lp{8cm}}
\hline
\textbf{Source} & \textbf{Data Provided} \\
\hline
DeFiLlama & TVL and APY for each token \\
Dune Analytics & Historical prices, protocol events, whale identification via SQL queries \\
Ethereum RPC & EOA verification via \texttt{eth\_getCode} \\
GraphQL & On-chain position data queried directly from protocol subgraphs \\
\hline
\end{tabular}
\end{table}
All scrapers are implemented in Python using the \texttt{dune-client} library for Dune Analytics queries, \texttt{web3.py} for RPC calls, and the \texttt{gql} library for GraphQL subgraph queries. Output is stored as JSON (metadata) and CSV (prices, events, whale lists).